\section{Method}
\label{sec:method}

\subsection{Overview}
\label{sec:overview}

We consider the class-incremental learning (CIL) setting: a model learns from a sequence of tasks $\{1, 2, \dots, T\}$, where each task $t$ introduces a disjoint set of classes $\mathcal{C}_t$. During training on task $t$, only data $\mathcal{D}_t = \{(x_i, y_i) \mid y_i \in \mathcal{C}_t\}$ is accessible, and data from prior tasks cannot be stored. At test time after task $T$, the model must classify any sample into the full label space $\bigcup_{\tau=1}^{T} \mathcal{C}_\tau$ without knowledge of task identity.

SPIE (Scalable Posterior Inference with Low-Rank Experts) addresses CIL through a three-stage pipeline applied at each incremental step.

At each task $t$, we first construct a lightweight expert $E_t$ by adapting a frozen pretrained ViT backbone through a compact, parameter-efficient module. This sidesteps full model expansion while preventing destructive overwriting of prior knowledge (Sec.~\ref{sec:experts}).

For each newly learned class, we extract a low-rank distributional prototype from the fixed backbone feature space and store it. These prototypes are replayed as synthetic features to keep the global classification boundary calibrated without retaining raw data (Sec.~\ref{sec:prototypes}).

At evaluation, SPIE estimates a soft task posterior rather than committing to a single hard task. It fuses all expert predictions in probability space, avoiding direct comparison of uncalibrated logits from different experts (Sec.~\ref{sec:posterior}).

\begin{figure}[t]
\centering
% \includegraphics[width=\linewidth]{figs/pipeline.pdf}
\caption{Overview of the SPIE framework. At each task $t$, (a) a per-task lightweight expert $E_t$ is constructed on a frozen backbone, (b) low-rank distributional class prototypes are extracted and stored, and (c) the global prototype bank is kept calibrated through prototype replay. At inference, a soft task posterior aggregates all expert predictions.}
\label{fig:pipeline}
\end{figure}

\subsection{Lightweight Stable-Plastic Experts}
\label{sec:experts}

SPIE addresses the stability-plasticity dilemma through two complementary components built on a shared frozen ViT backbone. A globally calibrated set of class-prototype vectors accumulates cross-task knowledge without parameter interference. These prototypes jointly serve as a stable classification boundary and the routing signal for inference. Per-task compact experts supply local discriminability by learning lightweight modulations of the frozen backbone weights.

The prototype bank consists of normalized weight vectors $w_j \in \mathbb{R}^d$, one per class, trained on the frozen backbone's CLS token features $z(x) = f_\theta(x) \in \mathbb{R}^{d}$. These vectors are cosine-normalized and serve as class-level representations after training and calibration. When task $t$ arrives with $|\mathcal{C}_t|$ new classes, we append $|\mathcal{C}_t|$ new weight rows and train them on $\mathcal{D}_t$ with a CosFace loss~\cite{wang2018cosface}:
\begin{equation}
\mathcal{L}_{\text{cos}} = -\log \frac{e^{s \cdot \cos(\theta_{y})}}{e^{s \cdot \cos(\theta_{y})} + \sum_{j \neq y} e^{s \cdot \cos(\theta_{j})}}, \quad \cos(\theta_j) = \langle \bar{z}(x), \bar{w}_j \rangle,
\label{eq:cosface}
\end{equation}
where $\bar{z}(x) = z(x)/\|z(x)\|_2$, $\bar{w}_j = w_j/\|w_j\|_2$, and $s$ is a learnable scale. After training, low-rank distributional prototypes of all seen classes are used to recalibrate the full prototype bank through distributional replay (Sec.~\ref{sec:prototypes}), correcting the bias toward recent tasks.

For the plastic path, each task $t$ receives a compact expert $E_t$ that adapts the frozen backbone through three coordinated lightweight mechanisms. The core of each expert is a low-rank adaptation module: for each MLP layer in the ViT backbone, we decompose its pretrained weight $W \in \mathbb{R}^{d_{\text{out}} \times d_{\text{in}}}$ into a frozen low-rank basis pair $B \in \mathbb{R}^{d_{\text{out}} \times r}$, $A \in \mathbb{R}^{r \times d_{\text{in}}}$ ($r=256$), and task-specific modulation is applied as
\begin{equation}
\Delta W^{(t)} = \big( \lambda_b^{(t)} \odot B \big) \cdot \text{diag}\big(\lambda_a^{(t)}\big) \cdot A,
\label{eq:adapt}
\end{equation}
where $\lambda_a^{(t)} \in \mathbb{R}^{r}$ and $\lambda_b^{(t)} \in \mathbb{R}^{d_{\text{out}}}$ are the only learnable parameters per layer. This constrains task adaptation to a compact subspace inherited from the pretrained weights, requiring merely $r + d_{\text{out}}$ scalars per layer. This yields approximately \textbf{3.3K} parameters per MLP block, compared to the millions required by full fine-tuning or standard LoRA adapters. In addition to the modulation scalars, $K=4$ learnable expert tokens $E^{(t)} \in \mathbb{R}^{K \times d}$ are appended to the input patch sequence. Through a structured attention mask (Figure~\ref{fig:attn_mask}), these tokens cross-attend to all backbone tokens to gather task-relevant information. Backbone tokens do not attend to expert tokens, preserving the integrity of the fixed feature space. The expert feature concatenates the backbone CLS token with the first expert token:
\begin{equation}
z_t^E(x) = \big[\, \text{CLS}(x) \;\big|\; E^{(t)}_1(x) \,\big] \in \mathbb{R}^{2d}.
\label{eq:expert_feat}
\end{equation}
A task-local classifier $\phi_t: \mathbb{R}^{2d} \rightarrow \mathbb{R}^{|\mathcal{C}_t|}$ then maps this expert feature to class logits within the current task's label space, trained with the same CosFace objective.

\begin{figure}[t]
\centering
% \includegraphics[width=0.85\linewidth]{figs/attn_mask.pdf}
\caption{Attention mask structure in SPIE experts. Backbone tokens attend to each other, expert tokens attend to both backbone tokens and other expert tokens, and backbone tokens \emph{do not} attend to expert tokens.}
\label{fig:attn_mask}
\end{figure}

At each incremental step, SPIE first learns the new class prototypes on $\mathcal{D}_t$ using CosFace training, then recalibrates the full prototype bank via distributional replay. Expert $E_t$ is subsequently trained on $\mathcal{D}_t$ with only its adaptation scalars, expert tokens, and local classifier $\phi_t$ updated, after which all its parameters are frozen. For a standard ViT-B/16 backbone ($d=768$, $d_{\text{mlp}}=3072$), each expert adds only \textbf{0.09\%} of the original backbone parameters. This enables scalable expansion without the linear parameter growth that burdens prior structural methods.

\subsection{Low-Rank Distributional Prototype Replay}
\label{sec:prototypes}

To keep the prototype bank calibrated across all seen classes without storing raw samples, SPIE maintains a compact distributional summary for each class. For each newly learned class $c$, we pass its training samples through the frozen backbone to collect CLS features $\{z_i\}_{i=1}^{n_c}$, and summarize the class as a Gaussian distribution $\mathcal{N}(\mu_c, \Sigma_c)$ with feature mean $\mu_c = \frac{1}{n_c}\sum_{i} z_i$.

Storing a full $d \times d$ covariance matrix is storage-prohibitive ($d=768$, $\sim$590K entries per class). We instead maintain a diagonal-plus-low-rank approximation. Let $Z_c$ denote the centered feature matrix. We compute its truncated SVD:
\begin{equation}
\Sigma_c \approx \text{diag}(\sigma_{\text{res}}^2) + V_r \cdot \text{diag}(\lambda_1, \dots, \lambda_r) \cdot V_r^\top,
\label{eq:lowrank_cov}
\end{equation}
where $V_r \in \mathbb{R}^{d \times r}$ contains the top-$r$ right singular vectors, $\lambda_i = s_i^2/(n_c-1)$ are the associated eigenvalues, and $\sigma_{\text{res}}^2$ is the residual diagonal variance. This representation requires merely $d \cdot (r+1)$ scalars per class ($r=16$, $\sim$13K), reducing storage by over \textbf{45$\times$} relative to the full covariance.

After learning task $t$, we sample synthetic features from the stored distributions of \emph{all} seen classes:
\begin{equation}
\tilde{z} = \mu_c + \varepsilon \odot \sigma_{\text{res}} \;+\; \eta \cdot \text{diag}(\sqrt{\lambda_1}, \dots, \sqrt{\lambda_r}) \cdot V_r^\top,
\label{eq:sampling}
\end{equation}
where $\varepsilon \sim \mathcal{N}(0, I_d)$ and $\eta \sim \mathcal{N}(0, I_r)$ are independent random vectors. We then fine-tune the prototype bank using standard cross-entropy loss on the synthetic feature set $\{(\tilde{z}_i, y_i)\}$ sampled uniformly across classes. This corrects the bias toward recent tasks without accessing past data, while the covariance structure captures class-conditional feature variation for more faithful replay. Importantly, these prototypes are \emph{not} used as nearest-prototype classifiers. They serve solely as distributional evidence for prototype calibration and, as we show next, as the foundation for expert-level posterior inference.

\subsection{Soft Posterior Inference}
\label{sec:posterior}

At evaluation, the model must classify a sample $x$ into $\bigcup_{\tau=1}^{T} \mathcal{C}_\tau$ without task identity. Directly combining uncalibrated logits across different experts is unreliable. Each local classifier $\phi_t$ is trained on its own task's class distribution and produces scores at incompatible scales. SPIE therefore avoids both hard task selection and naive logit aggregation. Instead, it derives a task-level routing signal from the calibrated prototype bank and lets expert predictions contribute in probability space, weighted by how confident the routing signal is for each task.

The calibrated prototype bank defines a unified feature space for task routing. Its weight vectors, after alignment, form a globally consistent class-prototype bank:
\begin{equation}
\mathcal{P} = \big\{ \bar{p}_c = w_c\,/\,\|w_c\|_2 \;\mid\; c \in \textstyle\bigcup_{\tau=1}^{T} \mathcal{C}_\tau \big\}.
\label{eq:protobank}
\end{equation}
For a test sample $x$, we compute its normalized shared feature $\bar{z}(x)$ and score every prototype via cosine similarity $a_c(x) = \langle \bar{z}(x), \bar{p}_c \rangle$. The routing signal for each task is the strongest class-level evidence within that task:
\begin{equation}
q(t \mid x) = \text{softmax}_t\!\left( \frac{s_t(x)}{\tau} \right), \quad s_t(x) = \max_{c \in \mathcal{C}_t} a_c(x),
\label{eq:task_posterior}
\end{equation}
where $\tau$ is a temperature parameter. This signal is not treated as a hard routing decision. When the prototype evidence strongly favors a particular task, $q(t \mid x)$ concentrates most weight on that task's expert. When the evidence is ambiguous for a sample that could plausibly belong to several tasks, the posterior remains diffuse, distributing weight across multiple candidates.

Given the expert logits $l^t(x) = \phi_t(z_t^E(x)) \in \mathbb{R}^{|\mathcal{C}_t|}$ from each task's expert, the final predictive distribution mixes expert-specific class probabilities weighted by the task posterior:
\begin{equation}
P(y \mid x) = \sum_{t=1}^{T} q(t \mid x) \cdot \frac{e^{l_y^t / \tau_e}}{\sum_{c \in \mathcal{C}_t} e^{l_c^t / \tau_e}} \cdot \mathbb{1}[y \in \mathcal{C}_t],
\label{eq:fusion}
\end{equation}
where $\tau_e$ controls the sharpness of per-expert predictions and $\mathbb{1}[y \in \mathcal{C}_t]$ restricts each expert to its own label space. This formulation naturally spans two regimes. When $q(t \mid x)$ is sharp, the prediction is dominated by a single expert. When $q(t \mid x)$ is diffuse, multiple experts jointly determine the outcome. The collective vote of several task-specialized classifiers effectively corrects an uncertain routing signal, preventing a single misrouting from catastrophically dominating the answer.

This design has three key properties. (i) Expert predictions are compared through calibrated probabilities, avoiding inter-expert scale discrepancies. (ii) The soft posterior adaptively blends routing confidence with expert specialization, so that hard-routing errors do not propagate. (iii) The routing signal $q(t \mid x)$ is derived from a single, globally calibrated prototype space $\mathcal{P}$, decoupling routing quality from expert training and avoiding the circular dependency of routing on uncalibrated expert outputs.
